%!TEX root = ./../../thesis-trier.tex

\todo[inline]{Hier an einigen Stellen nicht die Makros für semiautomaton und automaton benutzt.}
\IN{start state!of an automaton} 
\index{initial state|see{start state}}
\IN{final state set!of an automaton}
\IN{regular language}
\IN{accepted!language!by automaton}
\IN{accepted!word!by automaton}
\IN{automaton!as language acceptor}
\index{language acceptor automaton|see{automaton!as language acceptor}}
\index{recognized language|see{accepted language}}
\label{def:automaton-acceptor}
\Cref{intro:chap:short_intro_sync_aut} introduced semiautomata as mere finite-state transition systems.
Here we associate languages with automata.
A \emph{(deterministic and complete) finite automaton\footnote{In the literature, most of the time, automata are meant to be language
		acceptors without a special term and automata without
		a start state and a final state are called semiautomata.		
	} (as a language acceptor)} 
is a $5$-tuple $\automat{A}{Q}{\Sigma}{\delta}{q_0}{F}$ 
where $(Q, \Sigma, \delta)$ is a semiautomaton as defined before, $q_0$
is the \emph{start} (or \emph{initial) state} and $F \subseteq Q$ is the \emph{set of final states}. 
The language \emph{accepted}, or \emph{recognized}, by $\mathcal A$
is
\[
L(\mathcal A) = \{ u \in \Sigma^* \mid \applytrans{\delta}{q_0}{u} \in F \}, \label{def:accepted_lang_aut}
\]
i.e., the set of labels of paths going from the start state to some final state.
We also say a word $u \in \Sigma^*$ is \emph{accepted} by $\mathcal A$ if $u \in L(\mathcal A)$.\footnote{There should be no problem with using the modifier ``regular'' in the context of languages and in the context of permutation groups since both domains are sufficiently different
so that no confusion should arise.}	
A language $L \subseteq \Sigma^*$ is called \emph{regular} if there exists an automaton $\mathcal A$ 
such that $L = L(\mathcal A)$.\footnote{Eilenberg~\cite{Eilenberg1974} called
	languages acceptable by automata recognizable,
	which makes sense when considering more general monoids
	than the free monoid. Since we do not consider the more general case,
	we do not use this terminology here.	
	See also the related footnote on page~\pageref{def:RegSigma}.}	

Kleene~\cite{Kleene56} has shown that regular expressions denote precisely
the set of regular languages. A proof (in contemporary notation)
can be found in any textbook on formal language theory, e.g.~\cite{HopUll79,Pin2022}.

\begin{theorem}[Kleene's theorem~\cite{Kleene56}]
	% For every regular language $L \subseteq \Sigma^*$ there exists a regular
	% expression $R$ such that $
	The set of regular languages is precisely the set of languages
	describable by regular expressions, i.e., 
	for every $\automat{A}{Q}{\Sigma}{\delta}{q}{F}$
	there exists a regular expression $R$ such that $L(\mathcal A) = L(R)$
	and vice versa.
\end{theorem}


\IN{accessible!state!in automaton} 
\IN{accessible!automaton} 
\IN{coaccessible state!in automaton}
\index{initially connected|see{accessible automaton}}
Let $\automat{A}{Q}{\Sigma}{\delta}{q_0}{F}$
and $q \in Q$. The state $q$ is called \emph{accessible} if it is reachable from the start state, i.e.,  there exists a word $u \in \Sigma^*$
such that $\applytrans{\delta}{q_0}{u} = q$. 
It is called \emph{coaccessible} if there exists a word $u \in \Sigma^*$
such that $\applytrans{\delta}{q}{u} \in F$. We call $\mathcal A$ \emph{accessible} (or \emph{initially connected}) if all states are accessible. 
Obviously, for every regular language there exists an accessible automaton accepting it.


%Observe that a permutation automaton
%as a language acceptor is accessible if and only if it is strongly connected. \todo{im minimalaut teil sagen perm aut mit einem finalzsutand immer minimal}

For a semiautomaton $\semiautomat{A}{Q}{\Sigma}{\delta}$, $q \in Q$
and $S \subseteq Q$, by 
\[\mathcal A_{q,S} = (Q, \Sigma, \delta, q, S) \label{def:add_start_final}
\]
we denote the language acceptor automaton that results from $\mathcal A$ 
by choosing $q$ as a start state and $S$ as a set of final states.

There should be no confusion by calling both types of automata simply automata, as it should be clear
from the context whether we mean automata as language acceptors, or whether we are talking about 
semiautomata. Furthermore, both types are specified by a different sets of data, a $3$-tuple versus a $5$-tuple, which distinguishes them notation-wise. 
In case any ambiguity might arise, and it is necessary to differentiate between both notions,
we revert to using the proper names semiautomaton and automaton as language acceptor.
% or by specifying them explicitly  
%either as a $3$-tuple or a $5$-tuple. \todo{daraus kann man es ableiten.}
%However, if there might arise confusion because both notions appear in a single statement, 
%we use the more specific terminology.

Every automaton as a language acceptor is also a semiautomaton when forgetting about the start state
and the set of final states. So, all semiautomata notions transfer to automata as language acceptors
and we use them also for language acceptor automata in this sense.
%In this sense, all notions that we define and that do not involve a
%start state and a set of final state are to be regarded as definitions for both types of automata.\todo{english.}

Given automata $\automat{A_i}{Q_i}{\Sigma}{\delta_i}{q_i}{F_i}$
where $i \in \{1,\ldots,m\}$ and $m \ge 2$, 
for the product automaton $\mathcal A_1 \times \ldots \times \mathcal A_m$
we have the following:
\begingroup
\allowdisplaybreaks 
\begin{align}
	L(\mathcal A_1) \setminus L(\mathcal A_2) & = L((\mathcal A_1 \times \mathcal A_2)_{(q_1, q_2), ( F_1 \times Q )\setminus (Q \times F_2)}), \\
	\bigcap_{i=1}^m L(\mathcal A_i) & = L((\mathcal A_1 \times \cdots \times \mathcal A_m)_{(q_1,\ldots,q_m), F_1 \times \cdots \times F_m}), \label{prelim:eqn:product-intersection}	\\
	\bigcup_{i=1}^m L(\mathcal A_i) & = L((\mathcal A_1 \times \cdots \times \mathcal A_m)_{(q_1,\ldots,q_m), E}), 
\end{align}
\endgroup%
where $E = \bigcup_{i=1}^n Q \times \cdots \times \underbracket[0.4pt]{F_i}_{\clap{\scriptsize$i$-th position}} \times \cdots \times Q$.

Next, we state two results that relate the accepted language to morphic images and congruences.

\begin{proposition}
\label{prelim:prop:acceptor-morphism}
	Let $\varphi \colon Q \to S$ be a morphism between $\semiautomat{A}{Q}{\Sigma}{\delta}$
	and $\semiautomat{B}{S}{\Sigma}{\mu}$.
	Then $L(\mathcal A_{q_0,F}) \subseteq L(\mathcal B_{\applyfunc{q_0}{\varphi}, \applyfunc{F}{\varphi}})$
	for $q_0 \in Q$ and $F \subseteq Q$.
	Furthermore, if $\applyfunc{\applyfunc{F}{\varphi}}{\varphi^{-1}} = F$,
	then $L(\mathcal B_{\applyfunc{q_0}{\varphi}, \applyfunc{F}{\varphi}}) \subseteq L(\mathcal A_{q_0,F})$.
\end{proposition}
\begin{proof}
	Let $u \in \Sigma^*$.
	For $\applytrans{\delta}{q_0}{u} \in F$,
	we have $\applytrans{\mu}{\applyfunc{q_0}{\varphi}}{u} = \applyfunc{\applytrans{\delta}{q_0}{u}}{\varphi} \in \applyfunc{F}{\varphi}$ and so $u \in L(\mathcal B_{\applyfunc{q_0}{\varphi}, \applyfunc{F}{\varphi}})$.
	Next, if $\applyfunc{\applyfunc{F}{\varphi}}{\varphi^{-1}} = F$
	and $u \in L(\mathcal B_{\applyfunc{q_0}{\varphi}, \applyfunc{F}{\varphi}})$,
	then we get $\applyfunc{\applytrans{\delta}{q_0}{u}}{\varphi} = \applytrans{\mu}{\applyfunc{q_0}{\varphi}}{u} \in  \applyfunc{F}{\varphi}$
	and so $\applytrans{\delta}{q_0}{u} \in F$.
\end{proof}

\Cref{prelim:prop:congruence_induces_morphism} 
gives the following corollary to \cref{prelim:prop:acceptor-morphism}.

\begin{corollary}
\label{prelim:cor:acceptor-congruence}
  Let $\semiautomat{A}{Q}{\Sigma}{\delta}$
  be an automaton and $\sim$ be a $\TS{\mathcal A}$-congruence.
  Then for each $q \in Q$ and $F \subseteq Q$
  we have $L(\mathcal A_{q,F}) \subseteq L(( \mathcal A / {\sim} )_{[q]_{\sim}, [F]_{\sim}})$
  and if $F$ saturates~$\sim$, then
  $L(\mathcal A_{q,F}) = L(( \mathcal A / {\sim} )_{[q]_{\sim}, [F]_{\sim}})$ (note that $[F]_{\sim} = \{ [q]_{\sim} \mid q \in F\}$).  
\end{corollary}

Morphisms as introduced here are a concept defined for semiautomata. \Cref{prelim:prop:acceptor-morphism}
gives a way to extend this notions to morphisms for automata as language acceptors
by also stipulating that the start state is preserved
and the final state set saturates the kernel relation. We do not use such an extended notion of morphisms
for automata as language acceptors and thus can use the term morphism alone without any ambiguity here.